% Standalone problem document, generated from the adjacent README.md.
% Compile with XeLaTeX. Mathematical content is maintained in Markdown.
\documentclass[11pt,a4paper]{article}
\usepackage[fontsize=10.5pt]{scrextend}
\usepackage[margin=22mm,headheight=15pt,headsep=9mm,footskip=11mm]{geometry}
\usepackage{fontspec}
\setmainfont{texgyrepagella}[Extension=.otf,UprightFont=*-regular,BoldFont=*-bold,ItalicFont=*-italic,BoldItalicFont=*-bolditalic]
\setsansfont{texgyreheros}[Extension=.otf,UprightFont=*-regular,BoldFont=*-bold,ItalicFont=*-italic,BoldItalicFont=*-bolditalic]
\setmonofont{texgyrecursor}[Extension=.otf,UprightFont=*-regular,BoldFont=*-bold,ItalicFont=*-italic,BoldItalicFont=*-bolditalic,Scale=MatchLowercase]
\usepackage{amsmath,amssymb}
\usepackage{unicode-math}
\setmathfont{texgyrepagella-math.otf}
\usepackage{xcolor,microtype,enumitem,needspace,fancyhdr}
\usepackage{bookmark}
\definecolor{ink}{HTML}{17344B}
\definecolor{accent}{HTML}{197D80}
\definecolor{muted}{HTML}{596773}
\hypersetup{colorlinks=true,linkcolor=accent,urlcolor=accent,pdftitle={IV-04 - Exact
solution hulls for tridiagonal interval
systems},pdfauthor={Open Problems in Numerical Linear Algebra}}
\urlstyle{same}
\setlength{\parindent}{0pt}
\setlength{\parskip}{5pt plus 1pt minus 1pt}
\linespread{1.04}
\setlength{\emergencystretch}{3em}
\widowpenalty=10000
\clubpenalty=10000
\displaywidowpenalty=10000
\allowdisplaybreaks[1]
\setlist{nosep,leftmargin=1.5em}
\providecommand{\tightlist}{\setlength{\itemsep}{0pt}\setlength{\parskip}{0pt}}
\providecommand{\pandocbounded}[1]{#1}
\pagestyle{fancy}
\fancyhf{}
\fancyhead[L]{\sffamily\footnotesize\color{muted}OPEN PROBLEMS IN NUMERICAL LINEAR ALGEBRA}
\fancyhead[R]{\sffamily\bfseries\footnotesize\color{accent}IV-04}
\fancyfoot[L]{\parbox[b]{0.9\textwidth}{\sffamily\fontsize{8}{10}\selectfont\color{muted}Editorial ratings; a bounded search cannot certify that a problem remains unsolved.\\Literature check: 2026-09-10}}
\fancyfoot[R]{\sffamily\footnotesize\color{muted}\thepage}
\renewcommand{\headrulewidth}{0.3pt}
\renewcommand{\headrule}{\hbox to\headwidth{\color{accent}\leaders\hrule height \headrulewidth\hfill}}
\makeatletter
\renewcommand{\section}{\Needspace{5\baselineskip}\@startsection{section}{1}{0pt}{12pt plus 2pt minus 2pt}{4pt}{\sffamily\bfseries\large\color{ink}}}
\renewcommand{\subsection}{\Needspace{4\baselineskip}\@startsection{subsection}{2}{0pt}{9pt plus 1pt minus 1pt}{3pt}{\sffamily\bfseries\normalsize\color{ink}}}
\makeatother
\setcounter{secnumdepth}{0}
\begin{document}
{\sffamily\small\color{muted}Intervals and absolute value equations\par}
\vspace{3pt}
{\raggedright\hyphenpenalty=10000\exhyphenpenalty=10000\sffamily\bfseries\fontsize{21}{24}\selectfont\color{ink}Exact
solution hulls for tridiagonal interval systems\par}
\vspace{7pt}
{\sffamily\small\textbf{Difficulty:} challenging\quad\textbf{Importance:} interesting
to the community\par}
{\sffamily\small\bfseries\color{accent}Status: Partially resolved\par}
\vspace{4pt}
{\color{accent}\hrule height 0.8pt}
\vspace{6pt}
\textbf{Rating rationale:} Removing the remaining complexity obstruction
for genuinely tridiagonal interval systems appears to require new
analysis, warranting challenging. Community impact comes from exact
uncertainty bounds for a common structured linear system.

\textbf{Area:} interval linear systems; computational complexity

\subsection{Problem statement}\label{problem-statement}

The input is an \(n\times n\) tridiagonal interval matrix
\(\mathcal T\), \(n\ge1\), and an interval vector
\(\mathcal b\subset\mathbb R^n\), all with rational endpoints. Thus
\(T_{ij}=0\) for \(|i-j|>1\), and the remaining entries of \(T\) and all
entries of \(b\) range independently through their supplied closed
intervals. Define the united solution set

\[
\Sigma(\mathcal T,\mathcal b)=\{x\in\mathbb R^n:\exists T\in\mathcal T\ \exists b\in\mathcal b,\ Tx=b\}.
\]

Does a deterministic algorithm compute its exact coordinatewise interval
hull in time polynomial in the total binary input length? It must report
an empty solution set when appropriate; otherwise it must return

\[
\left[\inf_{x\in\Sigma}x_i,\ \sup_{x\in\Sigma}x_i\right],\qquad i=1,\ldots,n,
\]

with infinite endpoints explicitly represented. Finite endpoints are
returned exactly as rationals. No regularity promise is imposed:
intervals crossing zero and singular members are included. The
input/output convention makes the source's request for a polynomial
exact-hull algorithm precise; the target is the smallest box, rather
than an arbitrary enclosure.

This is distinct from computing the determinant range of the same
interval family. It asks for rigorous optimal error bars on the solution
vector of a structured uncertain linear system.

\newpage
\subsection{References}\label{references}

Jaroslav Horáček, Milan Hladík, and Michal Černý,
\href{https://arxiv.org/pdf/1602.00349}{\emph{Interval Linear Algebra
and Computational Complexity}}, §1.4.3, ``Structured systems,''
pp.~16--18; published in \emph{Applied and Computational Matrix
Analysis}, Springer (2017), 37--66,
\href{https://doi.org/10.1007/978-3-319-49984-0_3}{chapter DOI}.

Milan Hladík,
\href{https://doi.org/10.1007/978-3-030-31041-7_16}{\emph{An overview of
polynomially computable characteristics of special interval matrices}},
Springer (2020), 295--310;
\href{https://arxiv.org/pdf/1711.08732}{arXiv:1711.08732v1}, §2, p.~3,
second open question.

\subsection{Status check}\label{status-check}

Both sources explicitly leave the tridiagonal hull problem open;
bidiagonal systems are treated separately as polynomially solvable.
Searches on 2026-09-10 for tridiagonal interval systems, exact/tight
solution hulls, polynomial complexity, and 2025/2026 found no
classification. The existing IV-02 screen explains why later symbolic
algorithms using nonstandard interval multiplication do not settle
independent-entry range questions. General interval-system hardness
alone does not classify the tridiagonal restriction.

\subsection{Independent audit ---
2026-09-10}\label{independent-audit-2026-09-10}

The full 2017 chapter, §1.4.3, Theorem 9, proves exact hull computation
for regular bidiagonal intervals, a substantive subclass of this
tridiagonal target. Its proof uses division by diagonal intervals, and
the end of that section explicitly assumes bounded solution sets. The
treatment here of singular members and unbounded or empty solution sets
is an editorial extension of that chapter's question; the 2018
preprint's §2 states the question without an explicit regularity
promise. Even the general regular tridiagonal case remains unresolved in
the checked sources and subsequent searches.
\end{document}
